% !TEX root = ../Main.tex
\chapter{二次函数内接直线}

\ex{已知 $y_1=x^2$，$y_2=2x$，$y_1$ 与 $y_2$ 的交点记为 $A,B$。点 $P$ 在 $y_1$ 上、且在 $A,B$ 之间运动，求 $\triangle PAB$ 的面积的最大值。}

先定交点：$x^2=2x$ 给出 $x=0$ 或 $x=2$，故 $A(0,0)$，$B(2,4)$。

\begin{center}
\begin{tikzpicture}[>=Stealth,scale=0.95]
  \draw[->] (-0.6,0) -- (3,0) node[right]{$x$};
  \draw[->] (0,-0.5) -- (0,5) node[above]{$y$};
  \draw[thick,blue!70!black,domain=-0.55:2.45,smooth,samples=60] plot (\x,{\x*\x});
  \draw[thick,red!60!black,domain=-0.1:2.4] plot (\x,{2*\x});
  \coordinate (A) at (0,0);
  \coordinate (B) at (2,4);
  \coordinate (Q) at (1,2);   % on line y2
  \coordinate (P) at (1,1);   % on parabola y1
  \draw (A) -- (B);
  \draw (A) -- (P) -- (B);
  \draw[dashed] (P) -- (Q);
  \fill (A) circle (1.4pt) node[below left]{$A$};
  \fill (B) circle (1.4pt) node[above right]{$B(2,4)$};
  \fill (P) circle (1.4pt) node[below right]{$P$};
  \fill (Q) circle (1.4pt) node[above left]{$Q$};
\end{tikzpicture}
\end{center}

\sol{
    \textbf{法一（面积分割法）.} 过 $P$ 作竖直线交弦 $AB$ 于 $Q$。以竖直段 $PQ$ 为分割线，把 $\triangle PAB$ 沿 $PQ$ 分成左右两块，两块的水平"底宽"合起来正好是 $|x_B-x_A|$，于是
    \[
    S=\tfrac12\,|PQ|\cdot|x_B-x_A|.
    \]
    $|x_B-x_A|=2$ 是定值，所以\textbf{只需求出 $|PQ|$ 的最大值}。

    设 $x_P=p$（$0<p<2$）。$P$ 在抛物线上，$y_P=p^2$；$Q$ 在直线上，$y_Q=2p$。竖直段长
    \[
    |PQ|=y_Q-y_P=2p-p^2=-(p-1)^2+1.
    \]
    这是开口向下的抛物线，故 $|PQ|_{\max}=1$，在 $p=1$ 处取到。于是
    \[
    S_{\max}=\tfrac12\cdot 1\cdot 2=1.
    \]

    \textbf{法二（从两曲线的变化率来看）.} 记 $y_1=x^2$ 的变化率为 $k_1(x)=2x$，$y_2=2x$ 的变化率为 $k_2=2$（常数）。竖直间隙 $|PQ|=y_2-y_1$ 怎么变，取决于两条线谁涨得快：
    \begin{itemize}
        \item 若 $k_1(x)<2$，则 $y_1=x^2$ 涨得比 $y_2=2x$ 慢，间隙 $|PQ|$ 还在变大；
        \item 若 $k_1(x)>2$，则 $y_1$ 涨得更快，间隙 $|PQ|$ 开始缩小。
    \end{itemize}
    所以 $|PQ|$ 在 $k_1(x)=k_2$ 时取到最大。由 $2x=2$ 得 $x=1$，与法一的 $p=1$ 一致。
}
